\chapter*{Conclusion : La Liberté du "Non"}
\addcontentsline{toc}{chapter}{Conclusion : La Liberté du "Non"}

Au terme de ce voyage dans les profondeurs de nos neurones et le secret de nos histoires, une certitude demeure : l'affirmation de soi est un acte de d'amour. Amour de soi, pour ne plus se laisser piétiner, mais aussi amour de l'autre, pour lui offrir le cadeau de notre vérité.

Nous avons vu que le silence est une pathologie qui s'imprime dans nos chairs, sature nos hormones et éteint nos scans cérébraux. Mais nous avons aussi vu que la parole est une médecine. Chaque fois que Marc, Sophie ou Julie ont osé briser leur cage, ils n'ont pas seulement sauvé leur santé ; ils ont agrandi le monde.

La résilience n'est pas le retour à l'état antérieur. C'est l'invention d'un nouvel état, où la peur ne disparaît pas, mais où elle n'est plus la maîtresse de maison. S'affirmer, c'est choisir la vie contre l'inhibition. C'est accepter que le conflit soit parfois le prix de la rencontre. C'est, enfin, comprendre que notre dignité ne dépend pas de l'approbation des autres, mais de notre capacité à habiter pleinement notre existence.

\textit{Alors, n'ayez plus peur de votre propre voix. Elle est le souffle de votre liberté.}
